\section{有限序列}

\begin{definition}{有限序列,序列的长度}{}
	设$E$是集合,$\left(x_{i}\right)_{i\in I}$是$E$中的元素族.若$I$是由自然数构成的有限集,则称$\left(x_{i}\right)_{i\in I}$是$E$中的有限序列,$\left|I\right|$是该序列的长度.
\end{definition}

\begin{definition}{序列的项}{}
	设$\left(x_{i}\right)_{i\in I}$是集合$E$中长度为$n$的有限序列,$f:\left[1,n\right]\to I$是序同构.
	\begin{enumerate}
		\item 对任意$k\in\left[1,n\right]$,称$x_{f\left(k\right)}$是$\left(x_{i}\right)_{i\in I}$的第$k$项.
		\item $x_{f\left(1\right)}$和$x_{f\left(n\right)}$分别称为$\left(x_{i}\right)_{i\in I}$的首项和末项.
	\end{enumerate}
\end{definition}

\begin{convention}{}{}
	设$P\left\{i\right\}$是公式,则满足$P\left\{i\right\}$的元素$i$构成的集合$I$是$\N$的子集.我们把$E$中的序列$\left(x_{i}\right)_{i\in I}$写作$\left(x_{i}\right)_{P\left\{i\right\}}$,并称$x_{n}$是该序列的第$n$项.特别地,如果$I=\left\{n\in\N\middle|n\geq k\right\}$(其中$k\in\N$),则把$\left(x_{i}\right)_{i\in I}$记作$\left(x_{i}\right)_{n\geq k}$或$\left(x_{i}\right)_{n=k}^{\infty}$.
	
	同样的,对于集合族$\left(X_{i}\right)_{i\in I}$的乘积,相应地使用$\prod\limits_{P\left\{i\right\}}X_{i}$或$\prod\limits_{i=a}^{b}X_{i}$.类似的论断记号用法对$\left(X_{i}\right)_{i\in I}$的并,交,基数和,基数积以及代数学中的合成律都适用.
\end{convention}
























